% Generated draft: requires mathematical and source review.
\documentclass[11pt]{article}
\usepackage[margin=25mm]{geometry}
\usepackage{fontspec}
\setmainfont{Latin Modern Roman}
\usepackage{amsmath,amssymb,mathtools,unicode-math}
\setmathfont{Latin Modern Math}
\usepackage{enumitem,needspace,xurl,hyperref}
\setlength{\parindent}{0pt}
\setlength{\parskip}{5pt}
\setlength{\emergencystretch}{3em}
\newcommand{\N}{\mathbb{N}}
\newcommand{\Z}{\mathbb{Z}}
\newcommand{\Q}{\mathbb{Q}}
\newcommand{\R}{\mathbb{R}}
\newcommand{\C}{\mathbb{C}}
\newcommand{\F}{\mathbb{F}}
\newcommand{\A}{\mathbb{A}}
\newcommand{\PP}{\mathbb{P}}
\newcommand{\E}{\mathbb{E}}
\newcommand{\eps}{\varepsilon}
\newcommand{\dd}{\,\mathrm{d}}
\begin{document}
\section*{UnsolvedMath ID: 30004988}
\textbf{Source problem code:} catalogue entry
\subsection{Definitions and mathematical statement}
Let \(E\) and \(F\) be equivalence relations on sets \(M\) and \(N\), respectively. Write \([x]_E=\{x'\in M:x'Ex\}\) and \(M/E=\{[x]_E:x\in M\}\), and define \(N/F\) similarly. Once structures in a common signature have been specified on these quotients, an isomorphism \(\Phi:M/E\to N/F\) is a bijection preserving their specified operations, relations, and constants.

The supplied record calls these objects Borel quotient structures but does not specify their presentations, signature, admissible equivalence relations, or the meaning of topological triviality. These omissions prevent a definitive formal statement without additional source context.

A possible interpretation can be stated conditionally. Suppose that the intended presentations give \(M\) and \(N\) standard Borel structures, meaning measurable structures arising from Polish topologies. For an isomorphism \(\Phi:M/E\to N/F\), define its relation on representatives by
\[
\Gamma_{\Phi}=\{(x,y)\in M\times N:\Phi([x]_E)=[y]_F\}.
\]
Under the Borel-graph interpretation, \(\Phi\) is called topologically trivial when \(\Gamma_{\Phi}\) is Borel in the product standard Borel space \(M\times N\). This interpretation is a provisional formulation, not a definition verified from the supplied report passage. In particular, no additional Borelness assumption on \(E\) or \(F\) is being imposed here.

The source asks which assumptions ensure that every isomorphism between the intended quotient structures is topologically trivial. Under the provisional interpretation above, the target is to identify hypotheses \(H\), including the intended class of quotient presentations and any set-theoretic assumptions, such that for every admissible pair \((M,E),(N,F)\),
\[
H\ \Longrightarrow\ 
\bigl(\forall\Phi\in\operatorname{Iso}(M/E,N/F)\bigr)\,
\Gamma_{\Phi}\in\mathcal B(M\times N),
\]
where \(\mathcal B(M\times N)\) denotes the product Borel \(\sigma\)-algebra. The admissible class, the scope of \(H\), and the intended definition of topological triviality remain unresolved; this displayed template does not assert that any particular hypotheses suffice.
\subsection{Short English statement}
Determine assumptions ensuring that every isomorphism between the intended Borel quotient structures is topologically trivial. The source omits the precise class of quotient structures and the definition of topological triviality, so the Borel-graph formulation above is conditional and requires human mathematical review.
\subsection{Sources}
\begin{itemize}
\item Local UnsolvedMath registry, supplied complete record for permanent ID \(30004988\), titled \emph{Topological Triviality of Borel Quotient Isomorphisms}; exact source problem code \texttt{catalogue entry}. No individual proposer is identified; the record lists proposed year 2022.
\item Supplied report citation: \emph{Set Theory (2023)}, Oberwolfach Reports, DOI \texttt{10.4171/owr/2022/2}, pp.\ 79--106. \url{https://doi.org/10.4171/owr/2022/2}. The supplied background describes these reports as documenting mathematical workshops and their open-problem discussions. The title year 2023 and the record's proposed year 2022 are retained without reconciliation.
\item Supplied literature reference: Farah--Ghasemi--Vaccaro--Vignati, \emph{Corona Rigidity}. \url{https://arxiv.org/abs/2201.11618}.
\item Supplied literature reference: Farah--Shelah, \emph{Trivial automorphisms}. \url{https://arxiv.org/abs/1112.3571}.
\item Literature assessment dated exactly \texttt{2026-08-22}: the supplied curation labels the problem open and reports instances under forcing axioms and in important Boolean and \(C^*\)-quotient classes, while reporting no known general assumption-free characterization for arbitrary Borel quotient structures. This is supplied dated curation derived from the OWR record, not an independently verified assessment. The record expressly requests independent verification of time-sensitive status and cited claims.
\item The registry labels the statement \texttt{corrected\_verified} and reports direct checking of the official report, with definitions restored and extraction artifacts or source-dependent wording removed. Those are supplied verification claims; no source visit or independent verification is asserted here. The supplied excerpt itself does not contain the missing definitions.
\end{itemize}
\end{document}
